\section{模型检验与灵敏度分析}
\subsection{模型检验}
模型检验按“输出是否完整、约束是否满足、结论是否可复现”三项目标展开，分别对应三问的可核验要求。

问题一检验处理覆盖与时间结构。100条样本均形成输出，1920个已生成词区间全部满足非负起点、正长度、音轨范围内终点与起点随词序单调四项约束；典型样本给出词元、音频窗口与视频帧之间的映射，其字段与全量清单一致。区间生成率与结构约束共同刻画处理覆盖与一致性，6条超出接口长度样本的截断位置亦被逐条记录。

问题二同时报告分类与回归指标，并以配对方式核验差异。验证集Macro-F1为0.6012、强度MAE为0.5915；中性类F1为0.4848，低于负向与正向两类，221条涉及中性类的误判占全部误判的77.8\%，与错误归因中“零点附近判别更困难”的判断一致。补全、可靠度门控与蒸馏的作用通过缺失条件下的同掩码对照评价，详见第\ref{sec:q2}节。

问题三的检验独立于分类评价，用于确认解释是否忠实于预测行为。联盟价值分解的效率残差不超过$8.9\times10^{-16}$，满足Shapley效率性公理；以96条固定验证样本比较高分、随机与低分片段的累计删除效果，高分删除相对随机删除的AOPC差为0.1068，样本级95\%区间为[0.0887,0.1259]。条件充分性在最高档的绝对概率偏差为0.1608，说明主要模态中被舍弃的片段对预测仍有贡献；三种子与三种基线下的主参考模态一致率分别为82.29\%、83.33\%、86.46\%与100.00\%、90.62\%、87.50\%。

\subsection{灵敏度与稳健性分析}
灵敏度分析考察关键设置变化时结论是否保持稳定，包括采样强度、缺失强度与解释配置三个方面。

问题一在12条样本子集上增加边界扰动次数，视觉方差的排序相关达到0.995，质量分量的排序结构保持稳定，说明锚点间的相对比较不依赖扰动次数的具体取值。

问题二以两类干预刻画缺失强度的影响。连续缺口由1增至16个锚点时，F1由0.5906降至0.4462、MAE由0.6438升至0.8026；随机点状缺失比例由5\%增至40\%时，两项指标分别变化为0.6074至0.5469与0.6352至0.7316。两类干预均显示信息减少伴随整体性能退化，具体幅度随缺失模态与区间位置变化；在相近的实际新增遮挡量下，连续缺口的F1低于随机点状缺失，说明缺失在时间上的聚集方式与缺失数量共同影响预测，模型对分散缺失的容忍度高于连续缺口。

问题三在固定评估掩码下重复多种子训练与多基线解释，主参考模态的一致率保持在82\%以上，绝对贡献排序的Spearman相关为0.812与0.625，解释结论对训练随机性与基线选择均保持稳定。

\section{模型评价与推广}
\subsection{模型特点}
时间锚点模型以父词区间连接不同采样粒度，兼顾子词语义与原始音视频的追溯能力，覆盖与退化标记使局部观测缺失得到显式表达，而非被填充值掩盖。TSRF-D将观测节点补全、缺失状态编码与双任务学习结合，在保持输入接口一致的条件下完成局部缺失预测与全量专项推理，使鲁棒性设计不以常规性能为代价。ICF-Shapley模型以子集感知预测器统一分类与解释的输入状态，模态级贡献与片段级证据共享同一预测器，避免解释与预测脱节。三个模型依次解决时间组织、证据不完整与决策可追溯问题，共同构成从素材时间组织到不完整特征预测的处理框架。

\subsection{模型推广与改进}
三问共享的内容掩码、观测掩码与词级时间区间构成可直接复用的接口，模型可沿以下方向继续扩展。

在时间对齐方面，质量分数用于锚点间的相对排序，后续可引入数字读法归一化、长序列编码与人工词边界评价，扩大覆盖范围并给出绝对时间误差的标定。

在缺失预测方面，改进方向包括缺口上下文建模、中性类别边界刻画与门控收益评价，并可扩展至更长缺口、更多模态与多任务联合的预测场景；基于样本及多组掩码的不确定性分析，可区分缺失因素与采样波动的影响。

在解释方面，后续可增加随机排列重复、原词级片段对照与人工时间标注，把当前解释结果与人工判断对照，提升解释的稳定性与时间精度。
